% test-030.tex -- comandos \spcomb y \spfact
%
% Compilar:  lualatex-dev test-030.tex
% Validar:   show-pdf-tags --xml test-030.pdf | rnv-wrapp
%            verapdf --flavour ua2 --format text test-030.pdf
%
% Cada caso es un parrafo numerado, para poder referirse a el al
% escuchar el PDF. Los comentarios "esperado" son las lecturas previstas,
% aun NO verificadas con NVDA; reemplazarlas por "doc:" al escuchar.
% Las entradas invalidas (que detienen la compilacion) no van aqui.
\DocumentMetadata
  {
    lang=es-CL,
    pdfversion=2.0,
    pdfstandard=ua-2,
    tagging=on,
    tagging-setup={math/setup={mathml-SE}},
  }
\documentclass{article}
\usepackage[spanish]{babel}
\usepackage{unicode-math}
\usepackage{spintent}
\usepackage{hyperref}
\hypersetup
  {
    pdfdisplaydoctitle = true,
    pdftitle           = {test-030: combinatoria y factorial},
  }
\pagestyle{empty}
\begin{document}

\section*{A. Combinaciones}

% esperado (sin verificar en NVDA): «combinaciones de 5 en 2»
1. Combinación numérica (binomial): $\spcomb{5}{2}$.

% esperado (sin verificar en NVDA): «combinaciones de n en k»
2. Combinación con símbolos: $\spcomb{n}{k}$.

% esperado (sin verificar en NVDA): «combinaciones de n en k»
3. Estilo letra: $\spcomb[style=letter]{n}{k}$.

% esperado (sin verificar en NVDA): «combinaciones de 5 en 2»
4. Estilo binomial explícito: $\spcomb[style=binom]{5}{2}$.

% esperado (sin verificar en NVDA): «combinaciones con repetición de 5 en 2»
5. Combinación con repetición: $\spcomb[rep=true]{5}{2}$.

% esperado (sin verificar en NVDA): «combinaciones con repetición de 3 en 5»
6. Con repetición k puede superar a n: $\spcomb[rep=true]{3}{5}$.

\section*{B. Variaciones y arreglos}

% esperado (sin verificar en NVDA): «variaciones de 5 en 2»
7. Variación: $\spcomb[type=variation]{5}{2}$.

% esperado (sin verificar en NVDA): «variaciones con repetición de n en k»
8. Variación con repetición: $\spcomb[type=variation,rep=true]{n}{k}$.

% esperado (sin verificar en NVDA): «arreglos de 7 en 3»
9. Arreglo: $\spcomb[type=arrangement]{7}{3}$.

% esperado (sin verificar en NVDA): «arreglos con repetición de 7 en 3»
10. Arreglo con repetición: $\spcomb[type=arrangement,rep=true]{7}{3}$.

\section*{C. Permutaciones}

% esperado (sin verificar en NVDA): «permutaciones de n»
11. Permutación de n elementos: $\spcomb[type=permutation]{n}$.

% esperado (sin verificar en NVDA): «permutaciones de 5»
12. Permutación numérica: $\spcomb[type=permutation]{5}$.

% esperado (sin verificar en NVDA): «permutaciones de n en k»
13. Permutación de n tomando k: $\spcomb[type=permutation]{n}{k}$.

% esperado (sin verificar en NVDA): «permutaciones circulares de n»
14. Permutación circular: $\spcomb[type=permutation,circ=true]{n}$.

% esperado (sin verificar en NVDA): «permutaciones con repetición de 10 con repeticiones 3, 3 y 4»
15. Con repetición, lista numérica: $\spcomb[type=permutation,rep=true]{10}{3,3,4}$.

% esperado (sin verificar en NVDA): «permutaciones con repetición de n con repeticiones n sub 1 y n sub 2»
16. Con repetición, dos grupos: $\spcomb[type=permutation,rep=true]{n}{n_{1},n_{2}}$.

% esperado (sin verificar en NVDA): «permutaciones con repetición de n con repeticiones a, b y c»
17. Con repetición, tres símbolos: $\spcomb[type=permutation,rep=true]{n}{a,b,c}$.

\section*{D. Lectura y silencio}

% esperado (sin verificar en NVDA): «cinco sobre dos»
18. Lectura propia: $\spcomb[read-arg={cinco sobre dos}]{5}{2}$.

% esperado (sin verificar en NVDA): «(silencio)»
19. Silenciado: $\spcomb[silent=true]{5}{2}$.

\section*{E. Factorial}

% esperado (sin verificar en NVDA): «n factorial»
20. Factorial simbólico: $\spfact{n}$.

% esperado (sin verificar en NVDA): «5 factorial»
21. Factorial numérico: $\spfact{5}$.

% esperado (sin verificar en NVDA): «n menos 1 factorial (por verificar)»
22. Factorial de una expresión: $\spfact{(n-1)}$.

% esperado (sin verificar en NVDA): «(silencio)»
23. Factorial silenciado: $\spfact[silent=true]{5}$.

% esperado (sin verificar en NVDA): «n factorial»
24. Factorial con lectura propia: $\spfact[read-arg={n factorial}]{n-1}$.

\section*{F. Fórmulas}

% esperado (sin verificar en NVDA): «combinaciones de n en k es igual a n factorial entre k factorial por n menos k factorial (por verificar)»
25. Fórmula de la combinación: $\spcomb{n}{k} = \dfrac{\spfact{n}}{\spfact{k}\,\spfact{(n-k)}}$.

% esperado (sin verificar en NVDA): «probabilidad de A es igual a combinaciones de 5 en 2 entre combinaciones de 10 en 2 (por verificar)»
26. Probabilidad clásica: $\spProb{A} = \dfrac{\spcomb{5}{2}}{\spcomb{10}{2}}$.

% esperado (sin verificar en NVDA): «permutaciones de n es igual a n factorial (por verificar)»
27. Permutaciones y factorial: $\spcomb[type=permutation]{n} = \spfact{n}$.

\end{document}
